\section{Introduction}
\label{sec:intro}

Alman and Vassilevska Williams \cite{AV} gave the first truly subquadratic
algorithm for 3SUM and the first truly subcubic algorithm for APSP.  Both
follow from a deterministic algorithm for Exact Triangle in time
\(n^{3-\ETpaper}\) up to a logarithmic factor, which rests on a fast algorithm
for \emph{thin} matrix products: a recursion with Sch\"onhage's identity for
\(\schon313\oplus\schon141\), pruned to the entries that are wanted.  The
savings are small, and the paper does not optimize them.  The result has since
been formalized in Lean~4 as a set of statements about programs of a word RAM
\cite{upstream}.

This paper asks how far the same method goes with the same base identity.  It
has three ingredients, a determination of the optimal parameters, and the
all-edges route to APSP that footnote~10 of the paper mentions.

\subsection*{The ingredients}

\emph{An exact count of the leaves} (Section~\ref{sec:exact}).  The cost of the
recursion is governed by the numbers \(\beta_d\) of leaves of each order in a
tile.  The paper bounds their tail by a geometric series.  We bound it by its
exponential rate, a function \(\gtwo\) of the ratio \(c=L/m\) of the levels and
of the relative order (Theorem~\ref{thm:tail}); the proof is a comparison of
integers.  The paper's exponent is the tangent of the exact one at zero.  The
algorithm does not change: the same program gets a better bound on its running
time (Theorem~\ref{thm:offline}).

\emph{The true size of the middle part} (Section~\ref{sec:middle}).  The
reduction from Exact Triangle (Theorem~17 of the paper) produces instances
whose middle part has about \(D'/g'\) vertices, where \(D'\) is the dimension
parameter and \(g'\) the number of pieces per hash class, and calls the solver
as if it had \(D'\).  Passing the true size changes one assignment of the
program (Theorem~\ref{thm:mid}).  The solver then runs at a smaller inner
dimension \(D\) and, relative to it, with fewer wanted entries:
\(n^{2}/D^{\sigma}\) with \(\sigma>1/2\) instead of \(n^{2}/\sqrt D\).
Equivalently, the weights are hashed modulo a larger prime
(Remark~\ref{rem:prime}).

\emph{Pruned encodings} (Section~\ref{sec:pruned}).  The recursion reads the
encodings of the input only at the leaves with at most \(m\) symbols \(\Pz\).
They can be computed by a pruned form of Yates' algorithm in \(O(LM)\)
operations per band instead of \(O(10^{L})\) (Theorem~\ref{thm:pruned}), which
relaxes the condition under which a product counts as thin.  We write the
pruned encoder as a program and prove its running time on the machine
(Theorem~\ref{thm:encoder}); the rest of the solver is upstream's, re-verified
under the weaker contract that the encodings are right only at the leaves that
are read.

\emph{APSP through all-edges Exact Triangle} (Section~\ref{sec:alledges}).
The known reductions give the \(\minplus\)-product and APSP a third of the
saving of Exact Triangle.  Footnote~10 of the paper observes that the host of
its reduction also solves the all-edges version of Exact Triangle, to which the
\(\minplus\)-product reduces on the same number of vertices.  We write the two
programs (the all-edges host and the reduction, Theorems~\ref{thm:mid-ae}
and~\ref{thm:pairs}), and APSP keeps the whole saving.

\subsection*{The optimum}

With these ingredients the saving for Exact Triangle depends on four
parameters.  For a fixed ratio \(c\) its supremum has a closed form: it is
\(R(\Gam)\Gam/2\), where \(\Gam=\Gam(c)\) is the value at which two explicit
curves cross and \(R\) is the bound of the thinness condition
(Theorem~\ref{thm:sup}).  With pruned encodings the supremum over all ratios
lies in the interval \((\SstarLo,\SstarHi]\), and only ratios in
\((\cStarLo,\cStarHi)\) come within \(\SstarLo\)
(Theorem~\ref{thm:global}); with the paper's encodings it lies in
\((\SfullLo,\SfullHi]\), approached only for ratios in \((\cFullLo,\cFullHi)\)
(Theorem~\ref{thm:global-full}).  The theorems reach every saving below the
two suprema: the upstream program tests \(D^{18}\le N\) to decide whether its
recursion applies, which excludes the optimal ratios, and we replace the test
by one with a rational exponent (Section~\ref{sec:regime}).
So the method, with this identity and these ingredients, ends at a saving of
about \(\SprApprox\) for Exact Triangle, that is, at the exponent
\(2-\SprApprox/2\) for 3SUM.  Section~\ref{sec:limits} explains why no identity
of the same kind does better.

\subsection*{Results}

Table~\ref{tab:results} lists the exponents.  The column ``paper'' is
\cite{AV}; each of the next columns adds one ingredient to the previous one.

\begin{table}[ht]
\centering\small
\setlength{\tabcolsep}{5pt}
\begin{tabular}{lccccc}
\toprule
 & paper
 & Thm.~\ref{thm:T2}
 & Thm.~\ref{thm:T3}
 & Thm.~\ref{thm:T4}
 & Thm.~\ref{thm:T1}\\
\midrule
added & & parameters & exact count & middle size & pruned enc.\\
saving \(\delta\)
 & \(\ETpaper\) & \(\dTwo\) & \(\dThree\) & \(\SfullLo\) & \(\dOne\)\\
Exact Triangle
 & \(\ETexpPaper\) & \(\ETTwo\) & \(\ETThree\) & \(\ETFour\) & \(\ETOne\)\\
3SUM
 & \(\ThreeSumPaperStated\) & \(\ThreeSumTwo\) & \(\ThreeSumThree\)
 & \(\ThreeSumFour\) & \(\ThreeSumOne\)\\
\(\minplus\)-product, APSP
 & \(\APSPpaper\) & \(\APSPTwo\) & \(\APSPThree\) & \(\APSPFour\) & \(\APSPOne\)\\
\quad all-edges route
 & & & & \(\APSPFn\) & \(\APSPFnOne\)\\
\midrule
status & \cite{upstream} & checked & checked & checked & checked\\
\bottomrule
\end{tabular}
\caption{Exponents \(a\) of the running times \(O(n^{a})\) of deterministic
word-RAM programs.  The exponents of the column ``paper'' are those printed in
\cite{AV} and proved in \cite{upstream}.  The all-edges route
(Theorem~\ref{thm:apsp-ae}) is carried out for the two solvers of the last two
columns.}
\label{tab:results}
\end{table}

All four columns of new results hold without any hypothesis.  The first of
them uses none of the ingredients: it is the algorithm and the analysis of the
paper at better parameters.  The first three use the existing encoder, which is
the paper's; the last uses the pruned encoder of Theorem~\ref{thm:encoder}.
Theorems~\ref{thm:T4} and~\ref{thm:T1} give every saving below the two suprema
of Theorems~\ref{thm:global-full} and~\ref{thm:global}; the decimal exponents
are their values at one ratio each.

\subsection*{Formalization}

The results are formalized in Lean~4 \cite{Lean4} with Mathlib \cite{Mathlib},
on top of the formalization \cite{upstream} of the original paper; the
development and this paper are in the repository
\url{https://github.com/qawbecrdtey/ImprovedExponents}.  Every time
a result is machine-checked, the text names the Lean declarations in a
typewriter face, as in \lean{sum_beta_le_exp}.  The heading of a numbered
statement shows its status.
\begin{itemize}
\item \emph{No marker}: the statement is machine-checked.
\item ``\nf'': the statement is proved here on paper and is not
  machine-checked.  The marker never means ``unproved''.  It occurs only in
  Section~\ref{sec:limits}, for the mixtures of shapes, the hypothetical
  shapes, the experiments and a remark on \(\omega\); every numbered statement
  outside Section~\ref{sec:limits} is machine-checked, and so is the
  comparison of the shapes in Section~\ref{sec:limits}.
\end{itemize}
Section~\ref{sec:formalization} says what has to be trusted and how the
development is checked.

\subsection*{Organization}

Section~\ref{sec:framework} recalls the framework of \cite{AV}.
Sections~\ref{sec:exact}, \ref{sec:pruned} and~\ref{sec:middle} have the three
ingredients, Section~\ref{sec:optimum} the optimal parameters and
Section~\ref{sec:consequences} the theorems for 3SUM, Exact Triangle, the
\(\minplus\)-product and APSP, together with the all-edges route to APSP.
Section~\ref{sec:formalization} describes the formalization and
Section~\ref{sec:limits} the limits of the method.
